\chapter{Longitudinal interventions and repeated decisions}\label{ch:longitudinal}
\section{Why a sequence changes the problem}
In collections, yesterday's assistance affects today's arrears, today's arrears
influence the next assistance decision, and both predict eventual repayment.
Today's arrears are therefore both a consequence of earlier treatment and
a confounder for later treatment. Regressing repayment on every action and
every intermediate covariate can remove part of the total effect of earlier
assistance. Omitting the covariate can confound the later effect.

Let $L_t$ be information available immediately before decision $A_t$,
$H_t=(L_0,A_0,\ldots,A_{t-1},L_t)$, and let decisions occur at
$t=0,\ldots,T-1$. A policy $g_t(a\mid h)$ specifies a distribution over
feasible actions at each history. A deterministic rule is the special case
concentrated at $a=d_t(h)$. The target is
\begin{equation}
 \psi(g)=\E[Y^g],\qquad \Delta(g,g_0)=\psi(g)-\psi(g_0),
 \label{eq:policytarget}
\end{equation}
with population, horizon, outcome and eligibility stated. A comparison of
one action at baseline differs from a comparison of sustained strategies.

\section{Identification through the longitudinal g-formula}
Assume consistency, sequential exchangeability of actions given $H_t$,
and policy support: $g_t(a\mid h)>0$ only where the behavior probability
$e_t(a\mid h)>0$ at histories reached under $g$.
With no censoring for now, intervention replaces the assignment factors
in the observed law:
\begin{equation}
 p^g(l_0,a_0,\ldots,l_T)
 =p(l_0)\prod_{t=0}^{T-1}
   g_t(a_t\mid h_t)p(l_{t+1}\mid h_t,a_t).
 \label{eq:gformula}
\end{equation}
Integrating $Y$, included in terminal information $L_T$, under this law
gives $\psi(g)$ \citep{hernan2020}. The replacement is justified by causal
assumptions, not by a fitted transition model alone.

Two computational forms follow. G-computation fits outcome/covariate
transitions and simulates or integrates the intervention. Inverse weighting
uses the change of measure
\begin{equation}
 \psi(g)=\E\!\left[Y\prod_{t=0}^{T-1}
       \frac{g_t(A_t\mid H_t)}{e_t(A_t\mid H_t)}\right].
 \label{eq:longipw}
\end{equation}
The product cancels observed assignment factors, leaving the intervention
law. This derivation exposes both why weighting works and why it becomes
unstable when supported trajectories are rare.

Stabilization, clipping and normalization are different operations.
Clipping importance weights generally introduces bias relative to the
unclipped estimand; report the threshold and its effect. Effective sample
size $(\sum_iw_i)^2/\sum_iw_i^2$ describes weight concentration but is not
a replacement for a dependence-aware variance estimator or a proof of
identification.

\section{Recursive regression and orthogonal scores}
For a fixed policy $g$, define $V_T(H_T)=Y$ and work backward:
\begin{align}
 Q_t(h,a)&=\E[V_{t+1}(H_{t+1})\mid H_t=h,A_t=a],\\
 V_t(h)&=\sum_a g_t(a\mid h)Q_t(h,a),\qquad \psi(g)=\E[V_0(H_0)].
 \label{eq:recursiveq}
\end{align}
For continuous actions the sum becomes an integral. Define
$W_t^g=\prod_{s=0}^t g_s(A_s\mid H_s)/e_s(A_s\mid H_s)$.
A population orthogonal representation has score
\begin{equation}
 \Phi_g=V_0(H_0)+
 \sum_{t=0}^{T-1}W_t^g
       \{V_{t+1}(H_{t+1})-Q_t(H_t,A_t)\}.
 \label{eq:longdr}
\end{equation}
At true nuisance functions, each weighted residual has mean zero by
iterated expectation. Cross-fitted recursive estimators use this structure
to limit nuisance bias. Nested-function and longitudinal-rate analyses
explain the needed products of nuisance errors
\citep{nested2022,singh2022}. Sequential double robustness can hold when
an appropriate outcome or treatment component is correct at each stage
\citep{sequentialdr2018}; this property requires the corresponding recursive
construction and cannot be assigned to every algorithm labeled DR.

Longitudinal targeted estimation updates an initial outcome regression
along a fluctuation informed by treatment/censoring weights, aiming to
solve the influence-function estimating equation while respecting outcome
bounds. It is a useful alternative to a one-step estimator, especially
for bounded risks. The target, nuisance fits, targeting convergence and
support diagnostics must all be retained. A point estimate from a neural
sequence predictor does not supply these inferential properties.

Dynamic DML via g-estimation uses orthogonal residualization and recursive
removal of future effects in a structured dynamic model
\citep{dynamicdml2021}. Its independent-trajectory formulation suits some
many-customer panels. It does not directly supply a theorem for one
unrestricted aggregate bank series. The product must expose the chosen
effect parameterization, not relabel all dynamic models a single method.

\section{An exact two-decision example}
Set $L_0\sim\mathrm{Bernoulli}(1/2)$ and
\begin{align*}
 \Prob(A_0=1\mid L_0)&=0.2+0.6L_0,\\
 \Prob(L_1=1\mid L_0,A_0)&=0.1+0.5L_0+0.2A_0,\\
 \Prob(A_1=1\mid L_0,A_0,L_1)&=0.15+0.5L_1+0.1A_0,\\
 Y&=20+3A_0+4A_1+10L_1+6L_0.
\end{align*}
Every relevant action has positive probability. Under always versus never
treat, direct contributions add $3+4=7$ and the mediated change in $L_1$
adds $10(0.2)=2$. The total is $9$.
The intervention means are $35.5$ and $26.5$.

Enumerating all 16 observed cells verifies that the g-formula and exact
IPW reproduce both means. The observed always-adherer versus never-adherer
contrast is $\NumLongNaive$. A correctly specified conditional outcome
regression includes $L_1$ but its two action coefficients sum only to $7$;
that sum excludes the mediated component. The lesson concerns the causal
contrast, not an inability of regression to fit the observed outcome.

\begin{figure}[htbp]\centering
\includegraphics[width=\textwidth]{figures/effect_bias.pdf}
\caption{Synthetic static and longitudinal comparisons. Large apparent effects can reflect assignment and intermediate-variable conditioning.}
\label{fig:effects}
\end{figure}

\section{Feasible interventions and weak support}
An always-assist policy may be unsupported for low-risk customers or
ineligible accounts. Incremental interventions instead multiply treatment
odds by $\delta>0$:
\begin{equation}
 g_{\delta,t}(1\mid h)=
 \frac{\delta e_t(1\mid h)}{\delta e_t(1\mid h)+1-e_t(1\mid h)}.
 \label{eq:incremental}
\end{equation}
These preserve deterministic zero/one assignment regions
\citep{kennedy2019}. They answer how outcomes change under an odds shift,
not under universal treatment. Since the target itself depends on the
observed propensity law, its efficient score is not obtained by blindly
treating $g_\delta$ as externally fixed in Equation~\eqref{eq:longdr}.

Modified treatment policies can shift a rate or limit while respecting
conditional support \citep{lmtp2023}. A natural-value modification and
its stochastic analogue may require different sequential randomization
conditions. An economically feasible shift also needs a clear operational
implementation: a mathematical intervention on utilization is not necessarily
an action the bank can directly perform.

\section{Censoring, eligibility and competing events}
Account closure, debt sale, default, prepayment and loss of observation have
different meanings. If an outcome remains well-defined but becomes
unobserved, identification can require conditional independent censoring
and positive observation probability. Artificial censoring at deviation
from a treatment regime introduces an additional weighting or modeling task.
If treatment changes eligibility for a later action, simply restricting to
observed survivors changes the target population.

\citet{selective2025} explicitly model treatment-dependent eligibility;
their monotone eligibility setup must be checked against banking re-entry
and repeat episodes. \citet{competing2026} derive dynamic-regime cumulative
incidence methods with competing risks in a recent preprint. These motivate
separate outcome contracts for default and prepayment.

For disjoint terminal events with regime-specific marginal hazards
$\lambda^g_{D,k}$ and $\lambda^g_{P,k}$ among event-free accounts,
\begin{equation}
 F_D^g(K)=\sum_{k=1}^K
  \left[\prod_{j<k}(1-\lambda^g_{D,j}-\lambda^g_{P,j})\right]
  \lambda^g_{D,k}.
 \label{eq:cif}
\end{equation}
Estimating these regime hazards still requires longitudinal identification.
Censoring prepayment and calling the result the actual probability of default
would answer a different, potentially hypothetical question. A customer-value
outcome can incorporate both terminal events directly, provided cash-flow
ownership and observation are defined.

\section{Deep sequence models and the product boundary}
Counterfactual recurrent networks, causal transformers and neural
g-computation supply flexible representations or transition estimators
\citep{bica2020,melnychuk2022,gnet2020}. The inspected texts retain
consistency, sequential ignorability and overlap. Attention or adversarial
balance does not verify those assumptions.

They are appropriate challengers when long histories and nonlinear dynamics
defeat simpler models. Evaluate simulated interventions, misspecification,
support and policy value alongside factual prediction. Use simpler
g-computation, IPW and recursive DR comparators with the same target.
The first product should support a modest set of interpretable policies and
explicit horizons before claiming arbitrary long-horizon optimization.

\section{The estimating recursion and its error propagation}
For a cash-flow outcome, it is often more natural to keep rewards on their
actual dates than to aggregate them first into $Y$. Let $R_t$ be contribution
between decisions $t$ and $t+1$, let $h(H_T)$ settle the terminal account,
and fix discount factor $\gamma$. The action value satisfies
\[
 Q_t^g(h,a)=\E[R_t+\gamma V_{t+1}^g(H_{t+1})\mid h,a],
 \quad V_t^g(h)=\sum_a g_t(a\mid h)Q_t^g(h,a),
 \quad V_T^g=h.
\]
The terminal settlement must exclude cash already included in earlier
rewards. In lending, a remaining claim or expected recovery can belong
there; subtracting the same principal loss both monthly and at termination
would value a different account.

For fitted action values, define $\widehat V_t=\sum_a g_t\widehat Q_t$
and recursively set
\begin{equation}
 D_T=h,\qquad
 D_t=\widehat V_t(H_t)+\frac{g_t(A_t\mid H_t)}{\widehat e_t(A_t\mid H_t)}
       \{R_t+\gamma D_{t+1}-\widehat Q_t(H_t,A_t)\}.
 \label{eq:backward-dr}
\end{equation}
Expanding this recursion gives the full cumulative-ratio residual score.
It is the sequential forecast-and-correction construction of
\citet[Section 4.1]{jiang2016}, with the terminal account made explicit.
Every early ratio enters the correction for a later reward because early
actions determine which histories arrive at that later decision.

\begin{proposition}[Two routes to the mean of a fixed policy]
Under the stated sequential identification assumptions, with independent
evaluation trajectories, correct terminal settlement,
policy support and integrable terms, the score has mean $V^g$ if all
logging probabilities are correct. It also has mean $V^g$ if all action
values $\widehat Q_t=Q_t^g$ are correct and the fitted positive logging
probabilities are fixed history-measurable functions.
\end{proposition}
\begin{proof}
With correct ratios, each weighted dated quantity has its expectation under
the intervention law. Averaging $\widehat Q_t$ over the policy gives
$\widehat V_t$; the successive value terms telescope. With correct action
values, start at $T$ and use backward induction. The conditional expected
Bellman residual at each date is zero given $H_t,A_t$, so multiplying by
the history-measurable ratio leaves its mean zero.
\end{proof}

This is a mean identity. It does not state that every finite-sample normal
interval has nominal coverage, that a learned policy can be evaluated on
its own training outcomes, or that any estimated latent representation is
a valid state. The collections study will show visible undercoverage for
some unbiased or approximately unbiased constructions when path weights
are concentrated.

The propagation of nuisance error can be made explicit. Define
$b_t(h)=\E[D_t\mid H_t=h]-V_t^g(h)$ and
$\epsilon_t(h,a)=\widehat Q_t(h,a)-Q_t^g(h,a)$. Conditional on the
training fits,
\begin{align}
 b_t(h)&=\sum_a g_t(a\mid h)
       \left(1-\frac{e_t(a\mid h)}{\widehat e_t(a\mid h)}\right)
       \epsilon_t(h,a)\nonumber\\
 &\quad+\gamma\sum_a
   \frac{g_t(a\mid h)e_t(a\mid h)}{\widehat e_t(a\mid h)}
   \E[b_{t+1}(H_{t+1})\mid h,a],\qquad b_T=0.
 \label{eq:dynamic-remainder}
\end{align}
Substitution of $\E[D_{t+1}\mid H_{t+1}]=V_{t+1}^g+b_{t+1}$
into Equation~\eqref{eq:backward-dr} and collection of the regression
errors proves the identity. The first line is a local product of nuisance
errors; the second transports all later error backward. Very large early
ratios can amplify an otherwise modest continuation error. A small
one-step prediction loss is therefore an inadequate diagnostic for a
long-horizon causal value estimator.

\section{Learning the continuation rather than the whole trajectory}
Forward simulation of the g-formula needs a conditional distribution for
every covariate whose
future realization affects decisions or rewards. For a large mixed record
of balances, transactions, arrears and categorical eligibility, building a
coherent joint transition can be difficult. Recursive regression instead
learns only the conditional continuation needed for a fixed policy and
outcome. It can be much cheaper when the target is narrow. The price is
that a different policy, horizon or outcome can require a new sequence of
regressions.

A simulation model is more reusable for new policy scenarios but exposes
more mechanisms to misspecification. A reward model, balance evolution,
observation process and action eligibility must evolve together. Simulating
arrears independently of credit limits can generate impossible account
states. Neural g-computation offers flexible conditional laws
\citep{gnet2020}; flexibility does not remove the obligation to preserve
account identities and interventions. The product should compare simulated
policy values with recursive estimates when both are identified, not
judge the simulator only by factual one-step prediction.

For either approach, partition complete customer or household histories
before training. Within each training fold, generate future pseudo-outcomes
without using the held-out customer's outcomes. A naive implementation can
cross-fit the last regression while leaking the held-out outcome through
an earlier fitted continuation. The complete dependency graph of nuisance
fits must respect the fold boundary. Nested and longitudinal orthogonal
methods formalize rate and error-control requirements for such constructions
\citep{nested2022,singh2022}.

This distinction also affects hyperparameter tuning. Tuning on final policy
value can adapt to noise in the evaluation trajectories. Prediction-based
tuning within training folds preserves the declared evaluation, though it
may not minimize policy-value error. An outer evaluation of the entire
training and selection procedure is needed when the scientific claim is
about that procedure rather than one frozen policy.

\section{Targeted estimation and bounded outcomes}
For a binary terminal event, a fitted conditional probability should remain
between zero and one. A simple logistic fluctuation uses an initial
$\widehat Q_t^0$ and a one-dimensional parameter $\epsilon$:
\[
 \operatorname{logit}\widehat Q_t^\epsilon
   =\operatorname{logit}\widehat Q_t^0+\epsilon H_t^*,
\]
where the chosen covariate $H_t^*$ or equivalent observation weight is
constructed from the treatment and censoring mechanism for the target.
Fitting the fluctuation so the relevant weighted residual is zero aligns
the plug-in estimate with an influence-function equation. A longitudinal
procedure updates backward through the continuation regressions.

The displayed expression describes the mechanism, not a complete universal
TMLE algorithm. Different targets require different clever covariates,
weight conventions and recursion. The implementation must state whether
outcomes have been bounded or rescaled, what happens at fitted probabilities
zero and one, and how targeting convergence is checked. The relevant
sequential robustness analysis is estimator-specific
\citep{sequentialdr2018}. This is why a longitudinal adapter should retain
its native algorithm and diagnostics instead of wrapping every estimator
behind a function that returns only a scalar ``effect.''

Targeting can respect the parameter space while a one-step correction falls
outside it in a small sample. That is a useful finite-sample feature, but
does not repair practical positivity or unmeasured confounding. Conversely,
an out-of-range one-step estimate is a diagnostic of instability; silently
clipping it to a probability conceals the issue and changes inference.

\section{Observation loss has its own clock and model}
Let $C_{t+1}=1$ mean the next required record remains observed, conditional
on being observed through $t$. Write its conditional probability as
$c_t(H_t,A_t)$. If observation is conditionally independent of the relevant
future potential outcomes given that information, a change-of-measure
weight contains
\[
 W_t^{g,C}=\prod_{s=0}^{t}
 \frac{g_s(A_s\mid H_s)}{e_s(A_s\mid H_s)}
 \frac{C_{s+1}}{c_s(H_s,A_s)}.
\]
This formula separates two missing pieces: the action law and the reporting
law. Exact randomized treatment probabilities do not make a selectively
missing outcome ignorable. A customer whose debt is sold may have an
observed sale price, an unobserved later recovery and a changed economic
owner. The appropriate bank contribution can be observed at sale even
though a different lifetime recovery outcome is censored.

Administrative censoring at a fixed cutoff also changes maturity. A recent
cohort may have complete six-month outcomes but no twelve-month outcomes.
A twelve-month effect then needs older cohorts, a shorter target or an
extrapolation model. More recent transactions cannot fill calendar time
that has not elapsed. The data contract stores both the outcome end date
and the date by which the outcome could actually be known.

Competing events need the same discipline. Prepayment is not an arbitrary
missing default label. It changes the bank's remaining claim and may itself
be affected by assistance or pricing. Estimating default cumulative incidence,
event-free survival, a hypothetical no-prepayment risk and discounted bank
contribution are four different questions. Equation~\eqref{eq:cif} answers
the first under a regime; the cash-flow recursion can answer the fourth
without pretending that prepayment never occurs.

\section{Structural blips and dynamic effect parameters}
Sometimes the bank wants an interpretable effect parameter rather than a
complete regime value. Consider a two-action additive model in which the
conditional causal contribution of the final action is $\psi_1 A_1$.
Removing this contribution creates a blipped outcome
$Y^{\mathrm{blip}}=Y-\psi_1 A_1$. A g-estimation moment compares that
outcome across earlier actions after conditioning on their histories.
With an appropriate structural nested model, one estimates the final blip
first and moves backward, avoiding ordinary conditioning that blocks the
earlier mediated effect.

Dynamic DML combines this recursive removal with orthogonal nuisance
residualization \citep{dynamicdml2021}. The gains are a compact effect
parameterization and flexible adjustment; the costs are the restrictions
of the structural model, the need for support and dependence-compatible
sampling. Constant blips, heterogeneous blips and arbitrary dynamic regimes
are different model classes. The bank should choose among them according
to the decision and interpretability requirement rather than the apparent
modernity of the implementation.

The recommended comparison therefore has three tiers. An interpretable
structured model estimates a small set of economically meaningful dynamic
effects. Recursive DR or targeted estimation evaluates a prespecified set
of feasible regimes. A flexible transition or representation model explores
where those simpler constructions miss important history. Agreement is
useful evidence under shared assumptions; disagreement is a reason to
inspect targets, support and model restrictions, not to average incompatible
answers.
